\section{Architecture}
\label{sec:architecture}

The implementation is organized into five layers:

\begin{description}
  \item[\texttt{fields/}] Prime field $\mathbb{F}_p$ (radix-51
    Montgomery form) and quadratic extension $\mathbb{F}_{p^2}$.
  \item[\texttt{curves/}] Montgomery curves, $x$-only projective
    arithmetic, torsion bases, $2^e$-isogeny chains, and the
    \texttt{IsogenyDegree} newtype.
  \item[\texttt{surfaces/}] Abelian surfaces in theta coordinates,
    $(2,2)$-isogeny chains (Algorithm~8.47), gluing, and splitting.
  \item[\texttt{quaternions/}] Quaternion algebra $B_{p,\infty}$,
    lattices, ideals, HNF, L2 reduction, and \texttt{RepresentInteger}.
  \item[\texttt{deuring/}] The Deuring correspondence bridge:
    \texttt{IdealToIsogeny}, \texttt{FixedDegreeIsogeny},
    \texttt{SuitableIdeals}, action matrices.
\end{description}

\subsection{Randomness and the \texttt{\_derand} entry points}
\label{sec:randomness}

The SQIsign specification (Chapter~8) mandates
\textbf{AES256-CTR-DRBG} (NIST SP 800-90A, \S10.2.1) as the
pseudo-random number generator for optimized implementations.  The
C reference follows this verbatim: its NIST PQC test harness
(\texttt{PQCgenKAT\_sign.c}) instantiates CTR\_DRBG via
\texttt{randombytes\_init} with a 48-byte \texttt{entropy\_input}
and then services every call to \texttt{randombytes} from that
state.  Each KAT record in the \texttt{.rsp} files begins
\texttt{seed = <96 hex chars>} -- that 48-byte seed is the
CTR\_DRBG entropy input.

The 48 bytes are not arbitrary: for AES-256 the DRBG
\texttt{seedlen} equals $\text{keylen} + \text{blocklen} =
32 + 16 = 48$.  A shorter seed would leave part of the initial
CTR\_DRBG\_Update state at zero and diverge from the NIST harness.

We expose two parallel APIs for keygen and signing:
\begin{itemize}
  \item \texttt{SigningKey::generate<R: CryptoRngCore>(rng)} and
    \texttt{SigningKey::sign<R>(msg, rng)} sample 48 bytes of entropy
    from the caller's RNG and delegate to the derandomized entry
    points below.  Production callers pass \texttt{OsRng}.
  \item \texttt{SigningKey::generate\_derand(randomness: \&[u8;
    48])} and \texttt{SigningKey::sign\_derand(msg, randomness)}
    take the 48-byte seed directly, instantiate a local
    \texttt{Aes256CtrDrbg} from it, and drive all internal sampling
    from that DRBG.  These are the entry points that consume KAT
    seeds byte-for-byte.
\end{itemize}

Our \texttt{drbg} module implements AES256-CTR-DRBG directly on the
RustCrypto \texttt{aes} crate, which provides AES-NI intrinsics on
\texttt{x86\_64} and ARMv8 cryptographic-extension intrinsics on
\texttt{aarch64} out of the box.  The module is a direct
translation of the NIST reference (\texttt{CTR\_DRBG\_Instantiate},
\texttt{CTR\_DRBG\_Update}, \texttt{CTR\_DRBG\_Generate}): no
derivation function, no reseed counter, no additional input.  The
DRBG is exposed as an \texttt{rand\_core::RngCore}, so all
generic sampling code in \texttt{quaternions::lattice} (rejection
sampling, \texttt{rand\_interval} inside \texttt{reduce\_to\_prime\_norm},
\texttt{sample\_from\_ball}) drives it without further glue.
